% generated by scripts/make_example_figure.py -- do not edit by hand
\definecolor{probehelp}{HTML}{0072B2}
\definecolor{probeharm}{HTML}{D55E00}

\textbf{Question.} \emph{There are two warehouses. The first warehouse has twice as many boxes as the second warehouse. If the first warehouse has 400 boxes, how many boxes are there in both warehouses combined?}

\smallskip
\noindent\textbf{Rollout} (first 290 generated tokens; accepted
pivots in brackets, {\color{probehelp}blue} helpful and
{\color{probeharm}orange} harmful):

\smallskip
{\footnotesize\setlength{\parindent}{0pt}\textbf{\color{probeharm}[ Let]}'s see if I can figure this out. So, the problem says there are two warehouses. The first one has twice as many boxes as the second. And the first warehouse has 400 boxes. We need to find the total number of boxes in both warehouses combined. \\
\\
Okay, let me start by breaking down the information. Let's call the number of boxes in the second warehouse something. Maybe I can use a variable for that. Let me think... If the first warehouse has twice as many boxes as the second, then if I let the second warehouse be x boxes, the first one would be 2x. But wait, the problem says the first warehouse has 400 boxes. So that means 2x\textbf{\color{probehelp}[ equals]} 400.\textbf{\color{probehelp}[ Therefore]}, x would be 400 divided by 2, which is 200. So the second warehouse has 200 boxes. \\
\\
Then, to find the total number of boxes in both warehouses, I just add the boxes from the first and the second warehouses. So 400 plus 200. That equals 600. So the combined total is 600 boxes. \\
\\
Wait, let me check that again. If the first warehouse is twice the second, and the first is 400, then the second is half of 400, which is 200. Adding}

\smallskip
\begin{center}
\begin{tabular}{rlcccr}
\toprule
position & token & $p_{\text{before}}$ & $p_{\text{after}}$ & $\pdelta$ & rollouts \\
\midrule
39 & \texttt{ Let} & 0.42 & 0.05 & -0.38 & 17/40 $\to$ 2/40 \\
189 & \texttt{ equals} & 0.00 & 0.28 & +0.28 & 0/40 $\to$ 11/40 \\
195 & \texttt{ Therefore} & 0.40 & 0.62 & +0.22 & 16/40 $\to$ 25/40 \\
341 & \texttt{ =} & 1.00 & 0.00 & -1.00 & 40/40 $\to$ 0/40 \\
\bottomrule
\end{tabular}
\end{center}
